\documentclass[12pt,a4paper]{report}

% Fuentes CRÍTICAS para compatibilidad Word
\usepackage[utf8]{inputenc}
\usepackage[T1]{fontenc}
\usepackage{helvet}
\usepackage{courier}
\usepackage{mathptmx}
\usepackage{pifont}
\usepackage{fancyhdr}
\usepackage{mfirstuc}
\usepackage{datetime}

% Paquetes esenciales
\usepackage{tocbibind}
\DeclareRobustCommand{\mayusc}[1]{\MakeUppercase{#1}}
\DeclareRobustCommand{\titlecaps}[1]{\xcapitalisewords{#1}}
\usepackage{microtype}
\usepackage{geometry}
\usepackage{graphicx}
\usepackage{fancyhdr}
\usepackage{hyperref}
\usepackage{listings}
\usepackage{xcolor}
\usepackage{appendix}
\usepackage{float}
\usepackage{caption}
\usepackage{subcaption}
\usepackage{amsmath}
\usepackage{amssymb}
\usepackage{booktabs}
\usepackage{multirow}
\usepackage{enumitem}
\usepackage{titlesec}
\usepackage{wrapfig}
\usepackage{minitoc}
\usepackage{pdfpages}

% Configuración de márgenes
\geometry{
    left=2.5cm,
    right=2.5cm,
    top=2.5cm,
    bottom=2.5cm,
    headheight=15pt
}

% Configuración listings para código
\lstset{
    basicstyle=\ttfamily\fontsize{9}{11}\selectfont,
    backgroundcolor=\color{gray!10},
    frame=single,
    frameround=tttt,
    rulecolor=\color{gray!50},
    numbers=left,
    numberstyle=\tiny\color{gray},
    keywordstyle=\color{blue}\bfseries,
    commentstyle=\color{green!60!black}\itshape,
    stringstyle=\color{red},
    showstringspaces=false,
    breaklines=true,
    breakatwhitespace=true,
    tabsize=2,
    captionpos=b,
    xleftmargin=2em,
    framexleftmargin=1.5em,
    belowskip=1em,
    aboveskip=1em
}

% Headers y footers
\pagestyle{fancy}
\fancyhf{}
\fancyhead[L]{\color{blue!70!black}\textbf{\sffamily\leftmark}}
\fancyhead[R]{\color{blue!70!black}\textbf{\sffamily\thepage}}
\fancyfoot[C]{\color{gray}\small\sffamily wFabricSecurity - Zero Trust Security System}
renewcommand{\headrulewidth}{0.4pt}
renewcommand{\footrulewidth}{0.4pt}

% Estilos de capítulos y secciones
\titleformat{\chapter}[display]
{\normalfont\huge\bfseries\sffamily\color{blue!70!black}}
{\chaptertitlename\ \thechapter}{20pt}{\Huge}
\titlespacing*{\chapter}{0pt}{-30pt}{20pt}

\titleformat{\section}
{\normalfont\Large\bfseries\sffamily\color{blue!70!black}}
{\thesection}{1em}{}
\titlespacing*{\section}{0pt}{15pt}{10pt}

\titleformat{\subsection}
{\normalfont\large\bfseries\sffamily\color{orange!80!black}}
{\thesubsection}{1em}{}
\titlespacing*{\subsection}{0pt}{10pt}{5pt}

% Configuración hyperref
\hypersetup{
    colorlinks=true,
    linkcolor=blue!70!black,
    filecolor=blue!70!black,
    urlcolor=blue!70!black,
    citecolor=orange!80!black,
    pdftitle={wFabricSecurity - Zero Trust Security System for Hyperledger Fabric},
    pdfauthor={William Rodriguez},
    pdfsubject={Zero Trust Security, Hyperledger Fabric, Cryptographic Security},
    pdfkeywords={Zero Trust, Hyperledger Fabric, ECDSA, Cryptography, Security, Blockchain}
}

% Comandos personalizados
\newcommand{\code}[1]{\texttt{\fontfamily{pcr}\selectfont #1}}
\newcommand{\file}[1]{\texttt{\fontfamily{pcr}\selectfont #1}}
\newcommand{\pkg}[1]{\textsf{#1}}
\newcommand{\class}[1]{\textsf{#1}}
\newcommand{\method}[1]{\texttt{#1}}
\newenvironment{config}{\begin{table}[H]\centering\begin{tabular}{p{0.9\linewidth}}}{\end{tabular}\end{table}}

% Colores personalizados
\definecolor{primary}{RGB}{102,126,234}
\definecolor{secondary}{RGB}{118,75,162}
\definecolor{success}{RGB}{17,153,142}
\definecolor{danger}{RGB}{235,67,73}
\definecolor{warning}{RGB}{243,156,18}

\begin{document}

% Página de título
\begin{titlepage}
    \centering
    \vspace*{2cm}
    
    {\color{primary}\Huge\textbf{wFabricSecurity}}\\[0.5cm]
    {\color{secondary}\Large\textbf{Zero Trust Security System}}\\[1cm]
    {\Large\textit{for Hyperledger Fabric}}\\[2cm]
    
    \vspace{1cm}
    {\Large\textbf{Technical Documentation}}\\[0.5cm]
    {\large\textit{Complete Implementation Guide}}\\[2cm]
    
    \vspace{2cm}
    {\color{blue!70!black}\Large\textbf{Version 1.0.0}}\\[0.5cm]
    {\large\textbf{\today}}\\[3cm]
    
    \fbox{\begin{minipage}{0.8\linewidth}
        \centering
        \vspace{0.5cm}
        \Large\textbf{Cryptographic Identity \ding{51} Code Integrity \ding{51} Communication Permissions}\\
        \vspace{0.5cm}
        \end{minipage}
    }\\[2cm]
    
    \vfill
    \large\textbf{Maintained by: wisrovi}\\
    GitHub: github.com/wisrovi/wFabricSecurity
\end{titlepage}

% Página de autor
\begin{titlepage}
    \centering
    \vspace*{3cm}
    
    {\color{blue!70!black}\Large\textbf{Author}}\\[1.5cm]
    
    {\Huge\textbf{William Rodriguez}}\\[1cm]
    
    \textit{Lider de IA y Arquitecto de Soluciones}\\[0.3cm]
    \textit{eCaptureDtech}\\[0.2cm]
    \textit{Badajoz, Extremadura, Espana}\\[1.5cm]
    
    \begin{tabular}{c}
        \textbf{Contact Information:}\\[0.8cm]
        LinkedIn: es.linkedin.com/in/wisrovi-rodriguez\\[0.3cm]
        Email: william.rodriguez@ecapturedtech.com\\[0.3cm]
        GitHub: github.com/wisrovi\\[0.3cm]
        \href{https://github.com/wisrovi/wFabricSecurity}{Repository: github.com/wisrovi/wFabricSecurity}\\
    \end{tabular}\\[2cm]
    
    \vfill
    \fbox{\begin{minipage}{0.7\linewidth}
        \centering
        \Large\textbf{Expert in MLOps and Cloud Architecture}
    \end{minipage}
    }
\end{titlepage}

% Tabla de contenidos
\tableofcontents
\cleardoublepage

% Lista de figuras
\listoffigures
\cleardoublepage

% Lista de tablas
\listoftables
\cleardoublepage

% Lista de codigos
\lstlistoflistings
\cleardoublepage

% Agradecimientos
\chapter{Acknowledgments}
\cleardoublepage
\addcontentsline{toc}{chapter}{Acknowledgments}

\begin{center}
    \Large\textbf{Agradecimientos}\\[1cm]
\end{center}

Agradezco a todos los contribuidores del proyecto Hyperledger Fabric por crear una plataforma blockchain empresarial robusta. Tambien agradezco a la comunidad de codigo abierto por sus herramientas y librerias que hicieron posible este proyecto.

\bigskip

Este proyecto fue desarrollado como parte de un sistema de seguridad Zero Trust para entornos de fabricacion distribuida, con el objetivo de garantizar la integridad y autenticidad de todas las comunicaciones y operaciones.

\cleardoublepage

% Capitulo 1: Executive Summary
\chapter{Executive Summary}
\cleardoublepage

\section{Overview}
\index{Executive Summary}

\textbf{wFabricSecurity} es un sistema completo de seguridad Zero Trust dise\"nado para su implementacion con Hyperledger Fabric. Este sistema implementa validaciones criptograficas exhaustivas que garantizan la integridad del codigo, la autenticidad de las firmas digitales, los permisos de comunicacion y la verificacion de mensajes.

\section{Key Capabilities}

El sistema proporciona las siguientes capacidades fundamentales:

\begin{itemize}
    \item \textbf{Integridad de Codigo}: Verificacion SHA-256 del codigo fuente para detectar manipulacion maliciosa
    \item \textbf{Firmas ECDSA}: Criptografia de curva eliptica para firma y verificacion de mensajes
    \item \textbf{Permisos de Comunicacion}: Control de acceso granular (bidireccional, salida, entrada)
    \item \textbf{Integridad de Mensajes}: Verificacion de hash para detectar alteraciones en transmision
    \item \textbf{Rate Limiting}: Algoritmo de token bucket para proteccion contra ataques DoS
    \item \textbf{Retry Logic}: Reintentos automaticos con exponential backoff
    \item \textbf{Certificate Caching}: Cache LRU con TTL para optimizar el rendimiento
\end{itemize}

\section{Metrics}

\begin{table}[H]
    \centering
    \begin{tabular}{lr}
        \toprule
        \textbf{Metrica} & \textbf{Valor} \\
        \midrule
        Tests Totales & 228 \\
        Cobertura de Codigo & 84\% \\
        Calificacion Pylint & 9.29/10 \\
        Modulos Principales & 6 \\
        Excepciones de Seguridad & 8 \\
        \bottomrule
    \end{tabular}
    \caption{Metricas del Proyecto}
\end{table}

\cleardoublepage

% Capitulo 2: Introduction
\chapter{Introduction}
\cleardoublepage

\section{Background}

En entornos de fabricacion distribuida y sistemas empresariales complejos, la seguridad es una preocupacion critica. Los sistemas tradicionales confian automaticamente en las comunicaciones internas, lo que los hace vulnerables a ataques internos y compromisos de credenciales.

\textbf{Zero Trust} es un modelo de seguridad que elimina la confianza automatica. Cada peticion debe ser autenticada y autorizada antes de procesarse, independientemente de su origen.

\section{Objectives}

Este proyecto tiene como objetivos:

\begin{enumerate}
    \item Implementar un sistema de seguridad Zero Trust completo
    \item Proporcionar verificacion criptografica de codigo y mensajes
    \item Gestionar identidades y permisos de comunicacion
    \item Ofrecer proteccion contra ataques de denegacion de servicio
    \item Garantizar la trazabilidad mediante registro inmutable en blockchain
\end{enumerate}

\section{Scope}

El ambito de este documento incluye:

\begin{itemize}
    \item Arquitectura tecnica del sistema
    \item Guia de instalacion y configuracion
    \item Ejemplos de uso y casos de practica
    \item Referencia de API
    \item Mejores practicas y troubleshooting
\end{itemize}

\cleardoublepage

% Capitulo 3: Architecture
\chapter{Technical Architecture}
\cleardoublepage

\section{System Architecture}

El sistema sigue una arquitectura modular con las siguientes capas:

\begin{figure}[H]
    \centering
    \includegraphics[width=0.9\textwidth]{sources/architecture.pdf}
    \caption{Diagrama de Arquitectura del Sistema}
    \label{fig:architecture}
\end{figure}

\subsection{Presentation Layer}

Esta capa proporciona interfaces para interaccion con el sistema:

\begin{itemize}
    \item \class{CLI Tool}: Interfaz de linea de comandos
    \item \class{API Gateway}: Punto de entrada para aplicaciones externas
\end{itemize}

\subsection{Application Layer}

Contiene las clases principales del sistema:

\begin{itemize}
    \item \class{FabricSecurity}: Sistema completo Zero Trust
    \item \class{FabricSecuritySimple}: Version simplificada para casos basicos
\end{itemize}

\subsection{Security Layer}

Implementa todas las validaciones de seguridad:

\begin{itemize}
    \item \class{IntegrityVerifier}: Verificacion de integridad de codigo
    \item \class{PermissionManager}: Gestion de permisos de comunicacion
    \item \class{MessageManager}: Creacion y verificacion de mensajes
    \item \class{RateLimiter}: Limitacion de tasa de peticiones
    \item \class{RetryLogic}: Logica de reintentos automaticos
\end{itemize}

\subsection{Cryptographic Layer}

Servicios criptograficos fundamentales:

\begin{itemize}
    \item \class{HashingService}: Algoritmos de hash (SHA-256, BLAKE2)
    \item \class{SigningService}: Firmas ECDSA y HMAC
    \item \class{IdentityManager}: Gestion de certificados X.509
\end{itemize}

\section{Component Diagram}

\begin{figure}[H]
    \centering
    \begin{verbatim}
    +-------------------+
    |   CLI / API      |
    +--------+----------+
             |
    +--------v----------+
    |  FabricSecurity   |
    +--------+----------+
             |
    +--------v----------+
    | Security Layer    |
    | +-----------------+------------------+
    | Integrity | Permissions | RateLimiter  |
    +--------+----------+---------+---------+
             |
    +--------v----------+
    | Crypto Layer     |
    | Hashing | Signing | Identity
    +--------+----------+
             |
    +--------v----------+
    | Fabric Layer     |
    | Gateway | Network | Contract
    +--------+----------+
             |
    +--------v----------+
    | Storage Layer    |
    | Local | Fabric
    +-----------------+
    \end{verbatim}
    \caption{Diagrama de Componentes}
    \label{fig:components}
\end{figure}

\cleardoublepage

% Capitulo 4: Installation
\chapter{Installation Guide}
\cleardoublepage

\section{Requirements}

\subsection{Software Requirements}

\begin{itemize}
    \item Python 3.10 or higher
    \item pip (Python package manager)
    \item Docker and Docker Compose (optional, for Fabric network)
    \item Git
\end{itemize}

\subsection{Hardware Requirements}

\begin{itemize}
    \item CPU: 2 cores minimum
    \item RAM: 4 GB minimum
    \item Disk: 10 GB free space
\end{itemize}

\section{Installation Steps}

\subsection{1. Clone Repository}

\begin{lstlisting}[language=bash, caption={Clone Repository}]
git clone https://github.com/wisrovi/wFabricSecurity.git
cd wFabricSecurity
\end{lstlisting}

\subsection{2. Create Virtual Environment}

\begin{lstlisting}[language=bash, caption={Create Virtual Environment}]
python -m venv venv
source venv/bin/activate  # Linux/Mac
# venv\Scripts\activate   # Windows
\end{lstlisting}

\subsection{3. Install Package}

\begin{lstlisting}[language=bash, caption={Install Package}]
pip install -e .
\end{lstlisting}

\subsection{4. Install Development Dependencies}

\begin{lstlisting}[language=bash, caption={Install Dev Dependencies}]
pip install pylint black isort pytest pytest-cov
\end{lstlisting}

\section{Hyperledger Fabric Setup (Optional)}

Para ejecutar con la red blockchain de Fabric:

\begin{lstlisting}[language=bash, caption={Setup Fabric Network}]
cd enviroment
make setup    # Generate certificates
make up       # Start Docker network
cd ..
\end{lstlisting}

\cleardoublepage

% Capitulo 5: Usage Examples
\chapter{Usage Examples}
\cleardoublepage

\section{Basic Usage}

\subsection{Initialize Security System}

\begin{lstlisting}[language=python, caption={Initialize Security System}]
from wFabricSecurity import FabricSecurity

security = FabricSecurity(
    me="Master",
    msp_path="/path/to/msp"
)
\end{lstlisting}

\subsection{Register Identity and Code}

\begin{lstlisting}[language=python, caption={Register Identity and Code}]
# Register certificate
security.register_identity()

# Register code with hash
security.register_code(["master.py"], "1.0.0")
\end{lstlisting}

\subsection{Create and Verify Message}

\begin{lstlisting}[language=python, caption={Create and Verify Message}]
# Create signed message
message = security.create_message(
    recipient="CN=Slave",
    content='{"operation": "process_data"}'
)

# Verify message
if security.verify_message(message):
    print("Message is authentic and unaltered")
\end{lstlisting}

\section{Communication Permissions}

\begin{lstlisting}[language=python, caption={Communication Permissions}]
# Register permission
security.register_communication("CN=Master", "CN=Slave")

# Check permission
if security.can_communicate_with("CN=Master", "CN=Slave"):
    print("Communication allowed")
\end{lstlisting}

\section{Rate Limiting}

\begin{lstlisting}[language=python, caption={Rate Limiting}]
from wFabricSecurity import RateLimiter

limiter = RateLimiter(requests_per_second=100, burst=50)

# Acquire token (blocks if unavailable)
limiter.acquire()

# Try acquire (non-blocking)
if limiter.try_acquire():
    process_request()
\end{lstlisting}

\section{Retry with Exponential Backoff}

\begin{lstlisting}[language=python, caption={Retry Logic}]
from wFabricSecurity import with_retry

@with_retry(max_attempts=3, backoff_factor=2.0)
def unreliable_call():
    return fabric_invoke("RegisterTask", args)
\end{lstlisting}

\cleardoublepage

% Capitulo 6: API Reference
\chapter{API Reference}
\cleardoublepage

\section{Main Classes}

\subsection{FabricSecurity}

\class{FabricSecurity} es la clase principal que implementa el sistema completo de seguridad Zero Trust.

\textbf{Metodos Principales:}

\begin{itemize}
    \item \method{register\_identity()}: Registra el certificado de identidad
    \item \method{register\_code(paths, version)}: Registra el hash del codigo
    \item \method{verify\_code(paths)}: Verifica la integridad del codigo
    \item \method{create\_message(recipient, content)}: Crea un mensaje firmado
    \method{verify\_message(message)}: Verifica firma e integridad
    \item \method{can\_communicate\_with(sender, recipient)}: Verifica permisos
\end{itemize}

\section{Security Classes}

\subsection{RateLimiter}

Implementa el algoritmo de token bucket para limitacion de tasa.

\begin{lstlisting}[language=python, caption={RateLimiter Usage}]
limiter = RateLimiter(requests_per_second=100, burst=50)
limiter.acquire()
stats = limiter.get_stats()
\end{lstlisting}

\subsection{MessageManager}

Gestiona la creacion y verificacion de mensajes con TTL.

\begin{lstlisting}[language=python, caption={MessageManager Usage}]
manager = MessageManager(gateway, ttl_seconds=3600)
msg = manager.create_message(sender, recipient, content)
result = manager.verify_message(msg)
\end{lstlisting}

\section{Cryptographic Services}

\subsection{HashingService}

\begin{lstlisting}[language=python, caption={HashingService}]
service = HashingService()
hash_value = service.sha256("data")
message_hash = service.compute_message_hash("content")
code_hash = service.compute_code_hash(["file.py"])
\end{lstlisting}

\subsection{SigningService}

\begin{lstlisting}[language=python, caption={SigningService}]
service = SigningService(private_key=key)
signature = service.sign(data, signer_id)
is_valid = service.verify(data, signature, cert_getter, signer_id)
\end{lstlisting}

\cleardoublepage

% Capitulo 7: Best Practices
\chapter{Best Practices and Troubleshooting}
\cleardoublepage

\section{Best Practices}

\subsection{Security}

\begin{itemize}
    \item Siempre verificar el codigo antes de operaciones sensibles
    \item Mantener los certificados X.509 actualizados
    \item Implementar rotacion de claves periodica
    \item Usar HTTPS para todas las comunicaciones externas
\end{itemize}

\subsection{Performance}

\begin{itemize}
    \item Configurar el cache de certificados segun la carga esperada
    \item Ajustar los limites de rate limiting segun el tamano del sistema
    \item Implementar cleanup periodico de mensajes expirados
\end{itemize}

\section{Troubleshooting}

\subsection{Common Issues}

\begin{table}[H]
    \centering
    \begin{tabular}{p{0.3\linewidth}p{0.6\linewidth}}
        \toprule
        \textbf{Error} & \textbf{Solucion} \\
        \midrule
        CodeIntegrityError & El codigo fue modificado, verificar archivos \\
        PermissionDeniedError & Registrar permiso de comunicacion \\
        MessageIntegrityError & El mensaje fue alterado, retransmitir \\
        RateLimitError & Reducir frecuencia de peticiones \\
        \bottomrule
    \end{tabular}
    \caption{Problemas Comunes y Soluciones}
\end{table}

\section{Logging}

El sistema utiliza el modulo \pkg{logging} de Python. Configurar para produccion:

\begin{lstlisting}[language=python, caption={Logging Configuration}]
import logging

logging.basicConfig(
    level=logging.INFO,
    format='%(asctime)s - %(name)s - %(levelname)s - %(message)s'
)
\end{lstlisting}

\cleardoublepage

% Capitulo 8: Conclusions
\chapter{Conclusions and Future Work}
\cleardoublepage

\section{Conclusions}

Este proyecto implementa un sistema completo de seguridad Zero Trust para Hyperledger Fabric con las siguientes caracteristicas:

\begin{itemize}
    \item Verificacion criptografica de codigo y mensajes
    \item Gestion de identidades con certificados X.509
    \item Control de acceso granular mediante permisos
    \item Proteccion contra ataques DoS
    \item Registro inmutable en blockchain
    \item 228 tests automatizados con 84\% de cobertura
    \item Calificacion Pylint de 9.29/10
\end{itemize}

\section{Future Work}

Las siguientes mejoras estan planificadas para futuras versiones:

\begin{enumerate}
    \item Soporte para multiples organizaciones Fabric
    \item Implementacion de auditoria avanzada
    \item Integracion con sistemas SIEM
    \item Dashboard de monitoreo en tiempo real
    \item Soporte para hardware security modules (HSM)
\end{enumerate}

\cleardoublepage

% Bibliografía
\chapter{Bibliography}
\cleardoublepage
\addcontentsline{toc}{chapter}{Bibliography}

\begin{thebibliography}{99}

bibitem{fabric}
Hyperledger Fabric Documentation.
https://hyperledger-fabric.readthedocs.io/

bibitem{zerotrust}
NIST SP 800-207: Zero Trust Architecture.
National Institute of Standards and Technology.

bibitem{ecdsa}
Johnson, D., Menezes, A., \& Vanstone, S. (2001).
The Elliptic Curve Digital Signature Algorithm (ECDSA).
International Journal of Information Security.

bibitem{pep8}
PEP 8 -- Style Guide for Python Code.
https://www.python.org/dev/peps/pep-0008/

bibitem{pylint}
Pylint User Manual.
https://pylint.readthedocs.io/

bibitem{black}
Black - The Uncompromising Code Formatter.
https://black.readthedocs.io/

\end{thebibliography}

\cleardoublepage

% Apéndices
\begin{appendices}

\chapter{Configuration Reference}
\cleardoublepage

\section{Environment Variables}

\begin{table}[H]
    \centering
    \begin{tabular}{lp{0.6\linewidth}}
        \toprule
        \textbf{Variable} & \textbf{Descripcion} \\
        \midrule
        FABRIC\_PEER\_URL & URL del peer de Fabric \\
        FABRIC\_MSP\_PATH & Ruta al directorio MSP \\
        FABRIC\_CHANNEL & Nombre del canal \\
        FABRIC\_CHAINCODE & Nombre del chaincode \\
        RATE\_LIMIT\_RPS & Limite de peticiones por segundo \\
        RETRY\_MAX\_ATTEMPTS & Maximo numero de reintentos \\
        \bottomrule
    \end{tabular}
    \caption{Variables de Entorno}
\end{table}

\section{Config.yaml Template}

\begin{lstlisting}[language=yaml, caption={config.yaml}]
local_data_dir: /tmp/fabric_security_data
fabric_channel: mychannel
fabric_chaincode: tasks
fabric_peer_url: localhost:7051

retry_max_attempts: 3
retry_backoff_factor: 1.5
retry_initial_delay: 0.5

rate_limit_requests_per_second: 100
rate_limit_burst: 200

message_ttl_seconds: 3600

cert_cache_size: 100
cert_cache_ttl_seconds: 3600
\end{lstlisting}

\chapter{Exception Reference}
\cleardoublepage

\section{Security Exceptions}

\begin{itemize}
    \item \class{CodeIntegrityError}: El codigo fue modificado
    \item \class{PermissionDeniedError}: Sin permiso de comunicacion
    \item \class{MessageIntegrityError}: El mensaje fue alterado
    \item \class{SignatureError}: La firma es invalida
    \item \class{RateLimitError}: Limite de tasa excedido
    \item \class{RevocationError}: Participante revocado
    \item \class{ConfigurationError}: Error de configuracion
    \item \class{SecurityError}: Error general de seguridad
\end{itemize}

\chapter{Code Quality Metrics}
\cleardoublepage

\section{Coverage Report}

\begin{table}[H]
    \centering
    \begin{tabular}{lr}
        \toprule
        \textbf{Modulo} & \textbf{Cobertura} \\
        \midrule
        config & 98\% \\
        core & 91-100\% \\
        crypto & 72-91\% \\
        fabric & 63-78\% \\
        security & 64-88\% \\
        storage & 69-95\% \\
        \midrule
        \textbf{Total} & \textbf{84\%} \\
        \bottomrule
    \end{tabular}
    \caption{Resumen de Cobertura}
\end{table}

\section{Pylint Analysis}

La calificacion Pylint de 9.29/10 indica alta calidad de codigo con:

\begin{itemize}
    \item Nombres descriptivos y consistentes
    \item Documentacion completa de metodos
    \item Manejo adecuado de excepciones
    \item Estilo PEP 8 compliant (Black)
    \item Imports ordenados (isort)
\end{itemize}

\end{appendices}

\end{document}
